\begin{thebibliography}{99}
{\normalsize

\bibitem{ICT総研} ICT総研，2020年度 SNS利用動向に関する調査，https://ictr.co.jp/report/20200729.html，2021/01/03アクセス．

\bibitem{MeCab}MeCab: Yet another part-of-speech and morphological analyzer，https://taku910.github.io/mecab/， 2021/01/03 アクセス．

\bibitem{SlothLib}SlothLib，http://svn.sourceforge.jp/svnroot/slothlib/CSharp/Version1/SlothLib\\/NLP/Filter/StopWord/word/Japanese.txt， 2021/01/03 アクセス．

\bibitem{スポナビ}Sports navi，海外サッカートップ，https://soccer.yahoo.co.jp/ws/，2021/01/03アクセス．

\bibitem{Twitter}Twitter， Twitterについて， https://about.twitter.com/ja.html，2021/01/03アクセス．

\bibitem{Twitter API}Twitter Developer， Docs， https://developer.twitter.com/ja/docs，2021/01/03アクセス．

\bibitem{UEFA}UEFA，Country coefficients，\\https://www.uefa.com/memberassociations/uefarankings/country/\#/yr/2021，2021/01/03アクセス．

\bibitem{SVM} Bernhard E. Boser, Isabelle M. Guyon and Vladimir N. Vapnik，A training algorithm for optimal margin classifiers，In Proceedings of the 5th Annual Workshop on Computational Learning Theory，ACM Press，pp.142-152，1992．

\bibitem{ナイーブベイズ} George H. John and Pat Langley，Estimating continuous distributions in Bayesian classifiers，Proceedings of the Eleventh Conference on Uncertainty in Artificial Intelligence，pp.338-345，1995．

\bibitem{ランダムフォレスト} Leo Breiman，Random Forests，Machine Learning，45，pp.5–32，2001．

\bibitem{SMOTE} Nitesh V. Chawla，Kevin W. Bowyer，Lawrence O. Hall and W. Philip Kegelmeyer，SMOTE: Synthetic Minority Over-Sampling Technique，Journal of Artificial Intelligence Research，vol.16，pp.321-357，2002．

\bibitem{マクロ平均} Sebastian Raschka and Vahid Mirjalili，[第2版]Python機械学習プログラミング達人データサイエンティストによる理論と実践，株式会社インプレス，2018/3/21．

\bibitem{Jeon} Sungho Jeon，Sungchul Kim and Hwanjo，Spoiler Detection in TV Program Tweets，In: Information Sciences，vol. 329，pp.220-235，2016．

\bibitem{大島} 大島裕明， 中村聡史，田中克己，SlothLib: Webサーチ研究のためのプログラミングライブラリ，日本データベース学会 Letters，vol.6，no.1，pp.113-116，2007．

\bibitem{齊藤} 齊藤令，寺田実，Twitterの実況ツイートを利用したタイムライン上のネタバレ情報検知，情報処理学会第78回全国大会，pp.539-540，2016．

\bibitem{白鳥1} 白鳥裕士， 中村聡史，サッカーのネタバレが観戦者の態度に及ぼす影響，研究報告エンタテインメントコンピューティング(EC)，vol.2017-EC-43，no.17，pp.321-357，2017．

\bibitem{白鳥2} 白鳥裕士， 牧良樹， 阿部和樹， 中村聡史，ネタバレ確信度を考慮した試合実況データセット構築と分析手法の検討，ウェブインテリジェンスとインタラクション研究会予稿集，pp.1-6，2018．

\bibitem{中村} 中村聡史， 小松孝徳，スポーツの勝敗にまつわるネタバレ防止手法の検討: 情報曖昧化の可能性，情報処理学会論文誌，vol.54，no.4，pp.1402-1412，2013．
}

\end{thebibliography}